\begin{textsample}{POS Dim 2 – df – Score 20.00 – df/df\_em\_11.txt}  \label{ex:f2_pos_125}
2.6 EDUCAÇÃO \textbf{PROFISSIONAL} E TECNOLÓGICA NO CONTEXTO ATUAL Discutir os desafios atuais da educação, seja como \textbf{política} pública ou como processo de formação, em especial no que se refere aos jovens, exige necessariamente aprofundar o olhar sobre a Educação Profissional e Tecnológica ( EPT ) e sua organização no Brasil e no mundo.
Para tanto, cabe uma breve retomada da construção histórica dessa modalidade de ensino no país, que remonta ao período colonial, quando então o foco era voltado para as camadas mais pobres e para a mão de obra da produção manufatureira ( GARCIA, 2000, apud SILVA, 2018a ).
Em 1942, a \textbf{Reforma} Capanema promoveu a inclusão formal da EPT na estrutura de ensino do país, coincidindo, não por acaso, com a intensa industrialização ( CUNHA, 2014 ).
Nas décadas de 1980 e 1990, foram promovidas várias mudanças, as quais expandiram a EPT para estados e \textbf{municípios}.
A recente popularização dos \textbf{cursos} \textbf{técnicos} deu -se com a criação, em 2011, do Programa Nacional de Acesso ao Ensino \textbf{Técnico} e Emprego ( Pronatec ), que ampliou a \textbf{oferta} de \textbf{cursos} por meio das \textbf{redes} \textbf{federal} e \textbf{estaduais} e da parceria com o Sistema S ( FERRETI, 1997 ; BRASIL, 2011 ).
Apesar dos avanços robustos, em especial no que trata da \textbf{oferta} e do acesso aos \textbf{cursos}, a EPT ainda encontra dificuldades para se dissociar do forte apelo às camadas mais desfavorecidas que, em tese, têm mais dificuldade de acesso ao ensino superior.
Assim, a EPT se apresenta mais como um caminho terminativo do processo de formação, e não como uma etapa dessa formação ( SILVA, 2018a ).
A concepção utilitarista, focada apenas nos interesses do mercado de trabalho, encontra resistência de setores da sociedade e é combatida por diversos autores, com especial ênfase para Ciavatta ( 2005 ; 2011 ; 2014 ; 2019 ).
Essa linha de pensamento defende uma concepção política e ideológica mais ampla para a organização e o desenvolvimento da EPT, com base na educação integral e humanística, \textbf{alinhada} com as competências para o Século XXI, em seus grandes domínios, quais sejam, o cognitivo, o intrapessoal e o interpessoal ( UNESCO, 2015 ).
O potencial da EPT para o desenvolvimento socioeconômico pode ser verificado nas experiências de muitos países desenvolvidos, tais como a Alemanha, onde 54 \% da força de trabalho origina -se dessa modalidade, com \textbf{garantia} de formação integral e de \textbf{qualidade} ; e a Finlândia, que nos últimos 15 anos aumentou o quantitativo de 29 \% para 71 \% dos jovens entre 15 e 17 anos que \textbf{cursam} a EPT.
Apesar das dificuldades, a participação do Brasil na WorldSkills, maior competição mundial de EPT, realizada há 65 anos, rendeu aos brasileiros, selecionados pelo SENAI e SENAC, 136 medalhas em 18 participações, sendo que, na última edição, realizada em 2019, na Rússia, foram 63 jovens, competindo em 56 ocupações, das quais 13 ganharam medalhas e 28, certificados de \textbf{excelência} ( SENAI, 2020 ).
Nesse sentido, corrobora -se com a função educativa e social da EPT, buscando as condições para a formação cidadã, na qual os jovens possam se apropriar de todo o seu potencial, com especial recorte para as possibilidades de continuidade da formação, seja \textbf{aderindo} à \textbf{graduação} e aos seus desdobramentos, seja no pleno exercício da vida \textbf{profissional}, ou melhor ainda, em ambos os aspectos da vida.
Com o advento da Lei nº 13.415/2017, o maior desafio dos sistemas públicos de ensino recai na necessidade de organizar a \textbf{oferta} de EPT, de modo que a modalidade passe de fato a integrar a composição do Ensino Médio, deixando de ser mero complemento ou uma adição a ele, promovendo maior e melhor inclusão e \textbf{alinhada} às demandas atuais da sociedade ( AUR, 2018 ).
Entretanto, para que a integração do Ensino Médio com a EPT seja de fato estabelecida, a superação da histórica divisão entre a educação propedêutica e a educação \textbf{profissional} se impõe como condição imprescindível, ainda que a revisão das práticas empreendidas nesse sentido indique inúmeras dificuldades nas várias experimentações estudadas ( ARAÚJO ; FRIGOTTO, 2015 ).
O fato que precisa ser considerado de forma que a operacionalização da integração encontre condições para ser efetivada passa pelo que Costa ( 2012 ) reforça ao elencar que o ensino integrado precisa ir além das estratégias de organização de conteúdos escolares, avançando para a transformação da prática pedagógica, permeada pelos conteúdos éticos-políticos, ou seja : > [... ] a implementação do ensino médio integrado dentro de uma instituição não se resume à questão pedagógica, a um projeto curricular de ensino.
Requer a superação de diversos desafios, dentre eles os de gestão ; pedagógicos ; condições de ensino ; condições materiais ; hábitos estabelecidos culturalmente que limitam a formação integrada dos alunos ( COSTA, 2012, p. 38 ).
Assim, compreendendo as atuais demandas, a gestão da EPT vem atuando em parceria com as equipes do Ensino Médio e da Educação de Jovens e Adultos no sentido de encontrar caminhos para a superação de seus principais desafios institucionais, quais sejam : ampliação física, gestão de pessoas adequada às especificidades e demandas da EPT, contínua e permanente formação de docentes e gestores, \textbf{fortalecimento} dos processos pedagógicos e administrativos para \textbf{qualificação} do acesso e da formação, estabelecimento de parcerias sólidas com instituições públicas e privadas, institucionalização dos processos de \textbf{normatização} da EPT no âmbito do DF, alinhamento da \textbf{oferta} com a real demanda da sociedade e dos setores produtivos locais e regionais e permanente monitoramento da inserção dos egressos no mundo do trabalho ( SILVA, 2018b ).

% matched lemmas: aderir, alinhar, cursar, curso, estadual, excelência, federal, fortalecimento, garantia, graduação, município, normatização, oferta, política, profissional, qualidade, qualificação, rede, reforma, técnico
\end{textsample}
